% Standalone problem document, generated from the adjacent README.md.
% Compile with XeLaTeX. Mathematical content is maintained in Markdown.
\documentclass[11pt,a4paper]{article}
\usepackage[fontsize=10.5pt]{scrextend}
\usepackage[margin=22mm,headheight=15pt,headsep=9mm,footskip=11mm]{geometry}
\usepackage{fontspec}
\setmainfont{texgyrepagella}[Extension=.otf,UprightFont=*-regular,BoldFont=*-bold,ItalicFont=*-italic,BoldItalicFont=*-bolditalic]
\setsansfont{texgyreheros}[Extension=.otf,UprightFont=*-regular,BoldFont=*-bold,ItalicFont=*-italic,BoldItalicFont=*-bolditalic]
\setmonofont{texgyrecursor}[Extension=.otf,UprightFont=*-regular,BoldFont=*-bold,ItalicFont=*-italic,BoldItalicFont=*-bolditalic,Scale=MatchLowercase]
\usepackage{amsmath,amssymb}
\usepackage{unicode-math}
\setmathfont{texgyrepagella-math.otf}
\usepackage{xcolor,microtype,enumitem,needspace,fancyhdr}
\usepackage{bookmark}
\definecolor{ink}{HTML}{17344B}
\definecolor{accent}{HTML}{197D80}
\definecolor{muted}{HTML}{596773}
\hypersetup{colorlinks=true,linkcolor=accent,urlcolor=accent,pdftitle={TR-31 - The
four-exception conjecture for alternating tensor border-rank
varieties},pdfauthor={Open Problems in Numerical Linear Algebra}}
\urlstyle{same}
\setlength{\parindent}{0pt}
\setlength{\parskip}{5pt plus 1pt minus 1pt}
\linespread{1.04}
\setlength{\emergencystretch}{3em}
\widowpenalty=10000
\clubpenalty=10000
\displaywidowpenalty=10000
\allowdisplaybreaks[1]
\setlist{nosep,leftmargin=1.5em}
\providecommand{\tightlist}{\setlength{\itemsep}{0pt}\setlength{\parskip}{0pt}}
\providecommand{\pandocbounded}[1]{#1}
\pagestyle{fancy}
\fancyhf{}
\fancyhead[L]{\sffamily\footnotesize\color{muted}OPEN PROBLEMS IN NUMERICAL LINEAR ALGEBRA}
\fancyhead[R]{\sffamily\bfseries\footnotesize\color{accent}TR-31}
\fancyfoot[L]{\parbox[b]{0.9\textwidth}{\sffamily\fontsize{8}{10}\selectfont\color{muted}Editorial ratings; a bounded search cannot certify that a problem remains unsolved.\\Literature check: 2026-09-10}}
\fancyfoot[R]{\sffamily\footnotesize\color{muted}\thepage}
\renewcommand{\headrulewidth}{0.3pt}
\renewcommand{\headrule}{\hbox to\headwidth{\color{accent}\leaders\hrule height \headrulewidth\hfill}}
\makeatletter
\renewcommand{\section}{\Needspace{5\baselineskip}\@startsection{section}{1}{0pt}{12pt plus 2pt minus 2pt}{4pt}{\sffamily\bfseries\large\color{ink}}}
\renewcommand{\subsection}{\Needspace{4\baselineskip}\@startsection{subsection}{2}{0pt}{9pt plus 1pt minus 1pt}{3pt}{\sffamily\bfseries\normalsize\color{ink}}}
\makeatother
\setcounter{secnumdepth}{0}
\begin{document}
{\sffamily\small\color{muted}Tensor computations\par}
\vspace{3pt}
{\raggedright\hyphenpenalty=10000\exhyphenpenalty=10000\sffamily\bfseries\fontsize{21}{24}\selectfont\color{ink}The
four-exception conjecture for alternating tensor border-rank
varieties\par}
\vspace{7pt}
{\sffamily\small\textbf{Difficulty:} extreme\quad\textbf{Importance:} interesting
to specialist\par}
{\sffamily\small\bfseries\color{accent}Status: Partially resolved\par}
\vspace{4pt}
{\color{accent}\hrule height 0.8pt}
\vspace{6pt}
\textbf{Rating rationale:} Excluding all additional defective
Grassmannian secants is a longstanding classification barrier despite
extensive proved parameter ranges. Its direct importance is to
specialists in alternating tensor decompositions and secant geometry.

\subsection{Problem statement}\label{problem-statement}

For integers \(p\geq3\) and \(N\geq2p\), let \[
C_{p,N}=\{v_1\wedge\cdots\wedge v_p:
v_1,\ldots,v_p\in\mathbb C^N\}\subset\bigwedge^p\mathbb C^N.
\] For every integer \(s\geq1\), let
\(Y_s=\overline{\{x_1+\cdots+x_s:x_i\in C_{p,N}\}}^{\,\mathrm{Zar}}\).
Its expected affine dimension is \[
E_{p,N,s}=\min\left\{\binom Np,\ s\bigl(p(N-p)+1\bigr)\right\}.
\] Is \(\dim Y_s=E_{p,N,s}\) except precisely for the following four
triples, with the indicated actual affine dimensions?

\[
\begin{array}{c|r}
(p,N,s)&\dim Y_s\\\hline
(3,7,3)&34\\
(4,8,3)&50\\
(4,8,4)&64\\
(3,9,4)&74
\end{array}
\]

Equivalently, the listed cases should exhaust the defective secant
varieties of Grassmannians after using duality to impose \(N\geq2p\).
The restriction \(p\geq3\) excludes the separate, already understood
skew-symmetric matrix case. The dimension formula uses the affine cone,
hence the \(+1\) in each summand's parameter count.

\subsection{Why it matters}\label{why-it-matters}

This predicts the dimensions of spaces of alternating tensors that admit
a given number of decomposable wedge terms, including limiting
decompositions. It gives the alternating-tensor analogue of generic rank
formulas for symmetric tensors and identifies the exceptional formats
where parameter counting mispredicts approximation complexity.

\subsection{References and status
check}\label{references-and-status-check}

\begin{enumerate}
\def\labelenumi{\arabic{enumi}.}
\tightlist
\item
  K. Baur, J. Draisma, and W. A. de Graaf, \emph{Secant dimensions of
  minimal orbits: computations and conjectures},
  \href{https://doi.org/10.1080/10586458.2007.10128997}{Experimental
  Mathematics \textbf{16} (2007), 239--250}, Conjecture 4.1; the
  original statement uses \(p>2\) and \(2p\leq N\).
\item
  A. Oneto and E. Ventura, \emph{Ranks of tensors: geometry and
  applications}, \href{https://doi.org/10.1007/s40574-025-00472-9}{2025
  survey}, §2.2.3, Conjecture 2.16; its projective notation has
  \((k,n)=(p-1,N-1)\).
\item
  A. Taveira Blomenhofer and A. Casarotti, \emph{Nondefectivity of
  invariant secant varieties},
  \href{https://arxiv.org/html/2312.12335v2}{Full preprint, v2}, §4,
  Theorem 4.3 and Remark 4.4.
\end{enumerate}

\subsubsection{Status check ---
2026-09-10}\label{status-check-2026-09-10}

Checked the original-convention exception list against the 2025
restatement and Blomenhofer--Casarotti v2, Theorem 4.3 and Remark 4.4,
with targeted later Grassmannian-defectivity searches. Theorem 4.3
proves nondefectivity for substantive general rank ranges, so Partially
resolved applies to the displayed all-parameter target. These bounds
leave intermediate ranks and do not exhaust the possible exceptions. The
2025 survey still states the four-exception conjecture; no full
classification was located. The listed dimensions retain the affine-cone
convention.
\end{document}
